\chapter{Configuración y formato}
A continuación se detallan las características del formato del documento incluyendo la configuración de la página, márgenes, formato del texto, encabezados, y numeración. 

\section{Formato de página}

El tamaño del papel debe ser carta (21,6 cm x 27,9 cm), con fondo blanco. Las márgenes deben estar configuradas de la siguiente manera:

\begin{itemize}
    \item Superior: 2,5 cm
    \item Inferior: 2,5 cm
    \item Izquierdo: 2,5 cm
    \item Derecho: 2,5 cm
\end{itemize} 

\section{Formato del texto}
\begin{itemize}
    \item Tipo de fuente: Times New Roman, Open Sans o Verdana.
    \item Párrafos y textos en general (no títulos ni encabezados): fuente tamaño 12.
    \item Encabezado de capítulos: todo en mayúscula, tipografía en negrita tamaño 14.
    \item Encabezado de secciones principales dentro de un capítulo: palabras comienzan con mayúscula a excepción de artículos y preposiciones, tipografía en negrita de tamaño 14.
    \item Encabezado de secciones secundarias: palabras comienzan con mayúscula a excepción de artículos y preposiciones, tipografía en negrita de tamaño 12.
    \item Interlineado sencillo.
    \item Cada capítulo nuevo debe comenzar en una página nueva.
    \item El fin de una sección y el encabezado de la próxima no deben ser separados por espacios adicionales.
    \item Cada párrafo debe comenzar con una sangría de 1 cm a partir del margen izquierdo establecido.
    \item Cuando una sección comienza al final de una página, ésta debe ser trasladada a la página siguiente si el primer párrafo de la sección no alcanza a tener dos líneas de texto.
\end{itemize}

\section{Números de página}
Todas las páginas del documento deben estar numeradas, con excepción de la portada.  

La sección previa a la Introducción deberá numerarse con números romanos en minúscula. La numeración desde la introducción hasta el final del trabajo debe ser arábiga.  

Los números tanto romanos como arábicos deben estar en el pie de página, alineados abajo a la derecha. 

\section{Elementos del texto}
\subsection{Encabezados}
El trabajo debe estar dividido en secciones (o capítulos) y subsecciones. La subdivisión no debe pasar del tercer nivel. Las secciones principales, correspondientes a los capítulos, se numeran con números arábigos comenzando desde 1.

Las secciones resultantes de la división de los capítulos se enumeran secuencialmente con números enteros, comenzando desde 1 según su ubicación en el respectivo capítulo. Ejemplo: el capítulo 10 se subdivide en una sección secundaria 10.1, la cual se subdivide en 2 secciones terciarias 10.1.1. y 10.1.2.

\subsection{Ítems}
Dentro del texto, los ítems se pueden marcar de diferentes maneras. Por ejemplo: Utilizando números: 

1. Primer ítem. 

2. Segundo ítem. 

Utilizando letras: 

a. Primer ítem. 

b. Segundo ítem. 

Utilizando viñetas: 

Primer ítem. 

Segundo ítem. 

Es importante mantener un estilo de marcación consistente a lo largo del documento. La elección del sistema de marcación (números, letras o viñetas) debe basarse en la necesidad de destacar el orden o la jerarquía en la presentación de la información. 

\subsection{Notas al Pie de Página}
Las notas al pie de página son utilizadas para proporcionar aclaraciones o referencias adicionales sin interrumpir el flujo principal del texto. Se separan del cuerpo principal mediante una línea que ocupa aproximadamente un tercio del ancho de la página, partiendo del margen izquierdo. Estas notas se indican en el texto mediante superíndices. La tipografía de las notas al pie debe ser de tamaño 10. 

\subsection{Orientación de las Páginas}
La orientación normal de las páginas es vertical. Sin embargo, en casos excepcionales donde el contenido lo requiera, como en el caso de ciertas figuras o tablas, se puede usar una orientación horizontal. La numeración de las páginas debe seguir siendo secuencial y estar alineada abajo a la derecha en concordancia con el contenido de la página.

\subsection{Figuras y Tablas}
Son consideradas figuras: gráficos, diagramas, láminas, fotografías, esquemas de cualquier naturaleza, dibujos, planos, organigramas, flujogramas y cuadros tanto en color como blanco y negro.  

Tanto figuras como tablas deben incluir una leyenda y una numeración compuesta por el número del capítulo y el número de figura dentro del capítulo. La tipografía de la leyenda debe ser de tamaño 10. La leyenda debe localizarse abajo de la figura, justificada a la izquierda (en caso de que la leyenda sea de sólo una línea se puede colocar centrada). En tablas, la leyenda debe localizarse sobre la table, justificada a la izquierda. 

Tanto figuras como tablas deben ser ubicadas dentro del propio texto, lo más próximo posible del párrafo donde son referenciadas por primera vez. En el caso de que una figura o tabla no quepa completa en el espacio restante de una página, deberá desplazarse más adelante en el texto, de tal manera que el texto se acomode en dicho espacio restante y la figura pueda ubicarse en la página siguiente. 

El contenido de las figuras puede hacer uso de diferentes colores, y los textos que aparezcan deben presentar la misma nitidez que el texto principal del documento. 

Todas las figuras deben ser referenciadas en el texto principal. 

\subsection{Algoritmos}
En caso de que sea necesario presentar un pseudocódigo o algoritmo en el documento, éste deberá utilizar un tipo de letra de tamaño 12 y espaciado 1,15. Asimismo, el algoritmo debe estar indentado y se recomienda agregar los números de línea, de modo de facilitar su comprensión.  

Al igual que tablas y figuras, los algoritmos deben incluir una leyenda y una numeración compuesta por el número del capítulo y el número de figura dentro del capítulo. La tipografía de la leyenda debe ser de tamaño 12. La leyenda debe localizarse arriba del algoritmo, justificada a la izquierda. 
